\documentclass[11pt]{article}
\usepackage[margin=1in]{geometry}
\usepackage{amsmath,amssymb,amsthm,hyperref}
\newtheorem{theorem}{Theorem}
\newtheorem{lemma}[theorem]{Lemma}
\newtheorem{corollary}[theorem]{Corollary}
\newcommand{\E}{\mathbb E}
\newcommand{\Res}{\operatorname{Res}}
\begin{document}
\setlength{\emergencystretch}{1em}
\title{Strict Hankel positivity does not separate Appell moments}
\author{}\date{}\maketitle
\begin{abstract}
Consecutive Appell moment polynomials can share a nonreal quadratic
even when both their ordinary moment Hankel matrix and their
normalized-cumulant Hankel matrix are strictly totally positive.
Starting from a retained exact two-Gamma example, we construct a
countable positive Gamma convolution with this stronger property.
The positive weights, almost-sure convergence and moment generating
function are explicit. This refutes a strict-total-positivity
shortcut; it gives no single-Gamma-power or factorial SFC
counterexample.
\end{abstract}

For a positive random variable $Y$ with moments of every order, define
\[
 A_n(u)=\E(u+Y)^n,\qquad A_n'(u)=nA_{n-1}(u).
\]
When the moment generating function exists near zero, write
\[
 \log \E e^{zY}=\sum_{k\ge1}\kappa_k\frac{z^k}{k!},
 \qquad q_j=\frac{\kappa_{j+1}}{j!}\quad(j\ge0).
\]
Total nonnegativity of a Hankel matrix means nonnegativity of every
finite minor, with arbitrary strictly increasing row and column
indices. Strict total positivity means that every such minor is
positive.

\begin{lemma}[Positive-support Hankel minors]\label{lem:hankel}
Let $\rho$ be a finite positive measure on $(0,\infty)$ with every
moment finite. Its moment Hankel matrix is totally nonnegative.
Every minor of order $k$ is strictly positive if the support of
$\rho$ contains at least $k$ distinct points. In particular infinite
support implies strict total positivity at every order.
\end{lemma}
\begin{proof}
For increasing nonnegative integer lists $I=(i_1,\ldots,i_k)$ and
$J=(j_1,\ldots,j_k)$, determinant expansion and Fubini give
\[
 \det\left[\int x^{i_a+j_b}\,d\rho(x)\right]_{a,b=1}^k
 =\frac1{k!}\int
       \det[x_b^{i_a}]_{a,b=1}^k
       \det[x_b^{j_a}]_{a,b=1}^k\prod_{b=1}^k d\rho(x_b).
\]
Each generalized Vandermonde determinant is positive on
$0<x_1<\cdots<x_k$. Indeed a nonzero polynomial with at most $k$
nonzero terms has at most $k-1$ distinct positive zeros: divide by its
lowest monomial, differentiate, and apply Rolle's theorem inductively.
The determinant therefore never vanishes on this connected ordered
domain. Its sign is positive by coalescing the points and using
\[
 \det\left[
  \left(\frac{d}{dx}\right)^{b-1}x^{i_a}
 \right]_{a,b=1}^k
 =x^{\sum_a i_a-k(k-1)/2}
       \prod_{a<b}(i_b-i_a)>0.
\]
The product of the two determinants is nonnegative everywhere.
If there are $k$ distinct support points, disjoint ordered
neighborhoods of them have positive measure and give a strictly
positive contribution to the integral.
\end{proof}

\begin{theorem}[Recalled two-Gamma example]\label{thm:finite-obstruction}
Let $G_1,G_2$ be independent rate-one Gamma variables with shapes
$\alpha_1=125/9$, $\alpha_2=1/3249$, and set
\[
 Y=\tfrac12G_1+\tfrac{19}2G_2,\qquad
 \mu=\E Y=\frac{132}{19},\qquad x=u+\mu.
\]
The retained exact Gamma-convolution example gives
\[
 A_4(u)=(x^2-x+19)(x^2+x+3),\qquad
 A_5(u)=(x+2)(x^2-x+19)^2.
\]
Their exact monic gcd is $x^2-x+19$, with discriminant $-75$.
The ordinary moment Hankel matrix is strictly totally positive.
The normalized-cumulant Hankel matrix is totally nonnegative
and has rank two.
\end{theorem}
\begin{proof}
The parameters and factorizations are
Theorem \texttt{thm:gamma-convolution-obstruction} of the
\href{lower_exponent_initial_and_boundary.tex}{lower-exponent note}.
We recall its exact cumulants:
\[
 \kappa_1=\frac{132}{19},\quad
 \kappa_2=\frac72,\quad \kappa_3=4,\quad
 \kappa_4=\frac{81}4,\quad \kappa_5=582.
\]
They follow from
$\kappa_k=(k-1)!\sum_{i=1}^2\alpha_i\theta_i^k$ with
$\theta_1=1/2$, $\theta_2=19/2$.
The centered moments through degree five are
$1,0,7/2,4,57,722$, yielding the displayed factors.
The remaining factors are coprime to $x^2-x+19$.

The Gamma convolution has positive density on $(0,\infty)$, so
Lemma~\ref{lem:hankel} proves the ordinary moment assertion.
Also $q_j$ is the $j$th moment of the positive measure
$\sum_{i=1}^2\alpha_i\theta_i\delta_{\theta_i}$.
Its two distinct positive support points prove the cumulant
Hankel assertion by the same lemma.
\end{proof}

\begin{theorem}[Strict positivity of both Hankel matrices]
\label{thm:strict-obstruction}
There is a positive random variable $Y$, expressible as a countable
sum of independent scaled Gamma variables, whose moment generating
function exists near zero, such that both its ordinary moment Hankel
matrix and its normalized-cumulant Hankel matrix are strictly totally
positive, but $A_4,A_5$ have the common quadratic
\[
 (u+\E Y)^2-(u+\E Y)+19.
\]
\end{theorem}
\begin{proof}
We construct the required positive weights explicitly from four
moments, without assuming a truncated moment representation theorem.
Put
\[
 (m_0,m_1,m_2,m_3)
   =\left(\frac72,2,
                      \frac{27}{8},\frac{97}{4}\right).
\]
The three quantities
\[
 H_0=m_0m_2-m_1^2=\frac{125}{16},\qquad
 I_0=m_1m_3-m_2^2=\frac{2375}{64},\qquad
 J_0=m_0m_3-m_1m_2=\frac{625}{8}
\]
are positive. Define
\[
 \nu=\sum_{j\ge1}2^{-j}\delta_{1+1/j},\qquad
 r_k=\int t^k\,d\nu(t)\quad(0\le k\le3),\qquad
 \epsilon=\frac1{1000},\qquad M_k=m_k-\epsilon r_k.
\]
Here $0<r_k\le2^k$, so all $M_k$ are positive.
The quantities $H=M_0M_2-M_1^2$, $I=M_1M_3-M_2^2$ and
$J=M_0M_3-M_1M_2$ have the respective lower bounds
\[
 \begin{split}
 H&\ge H_0-\epsilon(4m_0+m_2)-4\epsilon^2>0,\\
 I&\ge I_0-\epsilon(8m_1+2m_3)-16\epsilon^2>0,\\
 J&\ge J_0-\epsilon(8m_0+m_3)-8\epsilon^2>0.
 \end{split}
\]
These follow by expanding the determinants and discarding favorable
positive terms.

Set $S=J/H>0$, $P=I/H>0$ and $p(t)=t^2-St+P$.
The definitions give
\[
 p(M_1/M_0)=-H/M_0^2<0.
\]
Thus $p$ has two distinct positive roots $t_1<t_2$ with
$t_1<M_1/M_0<t_2$. The weights
\[
 w_1=\frac{M_0t_2-M_1}{t_2-t_1},\qquad
 w_2=\frac{M_1-M_0t_1}{t_2-t_1}
\]
are positive. The measure
$\sigma_0=w_1\delta_{t_1}+w_2\delta_{t_2}$ has first two
moments $M_0,M_1$ by construction. The relation $p(t_i)=0$
and the definitions of $S,P$ show that its next two moments are
$M_2,M_3$. Hence
\[
 \sigma=\sigma_0+\epsilon\nu
\]
has first four moments $m_0,\ldots,m_3$ and infinite support in a
compact subset of $(0,\infty)$.

Enumerate its positive atoms as
$\sigma=\sum_\ell c_\ell\delta_{t_\ell}$, combining any duplicates.
Take independent
$H_\ell\sim\operatorname{Gamma}(c_\ell/t_\ell^2,1)$ and put
$Y=\sum_\ell t_\ell H_\ell$.
The sum is finite almost surely, since
$\E Y=\sum_\ell c_\ell/t_\ell<\infty$.
If $T=\sup_\ell t_\ell<\infty$, then for $|z|<1/(2T)$,
\[
 \sum_\ell \frac{c_\ell}{t_\ell^2}
                \bigl|\log(1-t_\ell z)\bigr|
 \le 2|z|\sum_\ell\frac{c_\ell}{t_\ell}<\infty .
\]
The logarithms therefore converge locally uniformly. For real
$0\le z<1/(2T)$, monotone convergence applied to the partial sums
identifies the resulting product with $\E e^{zY}$; it is finite.
The locally convergent logarithmic series gives, term by term,
\[
 \kappa_k=(k-1)!\int t^{k-2}\,d\sigma(t)\quad(k\ge1).
\]
In particular $\kappa_2,\ldots,\kappa_5$ equal those in
Theorem~\ref{thm:finite-obstruction}. Centering removes $\kappa_1$,
so the same factorizations of $A_4,A_5$ hold with $\mu=\E Y$.

The sequence $q_j$ is represented by $d\sigma(t)/t$, a finite
positive measure of infinite positive support. Lemma~\ref{lem:hankel}
proves strict total positivity of its Hankel matrix.
The law of $Y$ also has infinite positive support: condition on all
but one Gamma summand. Every open interval in $(0,\infty)$ has
positive probability for the whole sum. To see the needed small
values for the remaining sum, split it into finitely many summands
and a tail whose mean is less than a quarter of the desired upper
bound. Markov's inequality makes the tail smaller than half that
bound with positive probability; the finitely many summands are
simultaneously small with positive probability, independently.
The remaining Gamma summand has positive density on every positive
interval. A second application of Lemma~\ref{lem:hankel} proves the
ordinary moment assertion. Existence of the moment generating
function has already supplied every required finite moment.
\end{proof}

\begin{corollary}[The abstract shortcut is false]\label{cor:scope}
Even strict total positivity of both the ordinary moment Hankel
matrix and the normalized-cumulant Hankel matrix of a positive
random variable does not imply coprimality of its consecutive
Appell moment polynomials.
Neither example above is a single positive integer power $X^d$ of
one Gamma variable $X$. No counterexample to the binomial Gamma
family with shape $an+\beta$, or to factorial SFC, is asserted.
\end{corollary}
\begin{proof}
Theorem~\ref{thm:strict-obstruction} disproves the implication.
For $d\ge2$, $X^d$ has no finite positive exponential moment,
whereas both constructed variables do.
For $d=1$, a single scaled Gamma variable has a rank-one
normalized-cumulant Hankel matrix, whereas the first example has
rank two and the second has infinite rank.
The physical family with shape $an+\beta$ is therefore not
refuted by these constructions. They show that a proof of strict
positivity of its consecutive resultant must use additional
identities of that family.
\end{proof}

The exact replay checks the recalled rational parameters, the first five
cumulants, the two factorizations and gcd, the zero resultant, a
corrupted control, and the rational inequalities used in the infinite
construction. It does not replace the written convergence or
all-order Hankel proofs.

\end{document}
